% Generated by `taskgen course examples` from factoring.lenstra (seed 5).
% Do not edit: change examples.yaml and regenerate.
\begin{example}[алгоритм Ленстры]\label{ex:ecm-lenstra}
Найти нетривиальный делитель числа $N = 589$ алгоритмом Ленстры. Кривая $E_1\colon y^2 = x^3 + 3x + b_1$ над $\mathbb{Z}/N\mathbb{Z}$, где $b_1$ выбрано так, что точка $P_1 = (39, 378)$ лежит на $E_1$. Вычислять $Q_k = k\,Q_{k-1}$ ($Q_1 = P_1$, т.\,е. $Q_k = k!\,P_1$), $k = 2, 3, \dots$, удвоением-сложением по модулю $N$, пока не встретится необратимый знаменатель. Если алгоритм на $E_1$ не даст нетривиального делителя, взять кривую $E_2\colon y^2 = x^3 + 539x + b_2$ с точкой $P_2 = (365, 351)$.

\par\noindent\textit{Решение.} Формулы над $\mathbb{Z}/N\mathbb{Z}$: $\lambda = (3x_1^2 + a)/(2y_1)$ (удвоение), $\lambda = (y_2 - y_1)/(x_2 - x_1)$ (сложение), $x_3 = \lambda^2 - x_1 - x_2$, $y_3 = \lambda(x_1 - x_3) - y_1$. Деление на $d$ --- умножение на $d^{-1} \bmod N$ (расширенный алгоритм Евклида для $(N, d)$); если $\gcd(d, N) > 1$, обратного нет, и $\gcd(d, N)$ --- делитель $N$. Умножение на $k$ --- слева направо по двоичной записи $k$: $R \gets Q$, для каждого следующего бита $R \gets 2R$, при бите $1$ ещё $R \gets R + Q$.\par\medskip\noindent\textbf{Кривая $E_1$.} $b = y_0^2 - x_0^3 - a x_0 = 378^2 - 39^3 - 3 \cdot 39 \equiv 399\pmod{589}$, кривая $E_1\colon y^2 = x^3 + 3x + 399$; $4a^3 + 27b^2 \equiv 13$, $\gcd(13, 589) = 1$ (кривая невырождена по модулю обоих простых делителей). $Q_1 = P_1 = (39, 378)$; строки $k$ вычисляют $Q_k = k\,Q_{k-1}$, $Q = Q_{k-1}$, $m$ --- кратность точки $R = mP$ после шага, $d$ --- знаменатель:\[{\small\begin{array}{c|l|r|r|r|r|c}k & \text{шаг} & m & d & d^{-1} & \lambda & R = mP_1 \\ \hline 2 = 10_2 & 2R & 2 & 167 & 261 & 179 & (157, 293) \\ \hline 3 = 11_2 & 2R & 4 & 586 & 196 & 88 & (362, 515) \\  & R+Q & 6 & 384 & 385 & 524 & (172, 93) \\ \hline 4 = 100_2 & 2R & 12 & 186 & \gcd = 31 & \text{---} & \text{---} \\ \end{array}}\]\par Знаменатель $d = 186$ при $k = 4$ необратим: $\gcd(186, 589) = 31$ (остатки алгоритма Евклида: $589, 186, 31, 0$). Объяснение: $\#E(\mathbb{F}_{19}) = 20$, $\operatorname{ord}(P_1 \bmod 19) = 5$; $\#E(\mathbb{F}_{31}) = 24$, $\operatorname{ord}(P_1 \bmod 31) = 2^{2} \cdot 3$; на последнем шаге вычисляется точка $12P_1$ ($12 = 2^{2} \cdot 3$), и она равна $\mathcal{O}$ по модулю $31$, но не по модулю $19$.\par Делитель $31$, $N / 31 = 19$.\[N = 19 \cdot 31 = 589.\]
\end{example}
